\section{Introduction}
An acronym is a short written form that stands for a longer expression. In Hebrew, acronyms are written with a quotation mark before the last letter, as in \heb{ח״כ} (\heb{חבר כנסת}, ``member of parliament''), and they are common in news, parliamentary and everyday text. Many acronyms are ambiguous. For example, \heb{מט״ח} can mean \heb{מטבע חוץ} (foreign currency) or \heb{המרכז לטכנולוגיה חינוכית} (the educational technology center). Only the surrounding sentence tells the reader which one is meant.

Hebrew makes this problem harder than it is in many languages. Short words are written without vowels, so unrelated expansions often share the same letters. Prefixes such as ``in'', ``to'' and ``the'' attach directly to the word, so the acronym in a sentence may differ from its dictionary form. Hebrew also has far fewer resources than English, so a system cannot rely on large annotated corpora.

A language model can expand an acronym in two ways. In \emph{free generation}, it reads the sentence and writes the expansion. In \emph{candidate selection}, it is also given a list of possible expansions and picks one. Generation can produce an answer that is not on any list, but it can also produce a wrong or unsupported one. Selection limits the answer to the list, so it depends on the list being complete. Prior work has not directly compared the two formats on the same occurrences.

We run this comparison on a human-reviewed held-out test set of 381 Hebrew sentences. Most (77\%) occur naturally, in parliamentary proceedings and on Wikipedia, and the rest were written by an LLM to add a second sense to types that had only one. We test three open LLMs, one of them the Hebrew-specialised DictaLM, four proprietary LLMs, and a pre-trained Hebrew encoder, which we evaluate both before and after fine-tuning. All systems see the same sentence and the same marked acronym. No acronym form in the test set appears in training, so the encoder is tested on acronyms it has never seen.

Our findings are as follows. First, a small encoder fine-tuned on fewer than 3,000 examples reaches 71.7\% in candidate selection, up from 64.3\% before fine-tuning. That is higher than every open LLM, including the much larger Hebrew-specialised DictaLM 2.0 (58.5--69.8\%), though lower than every proprietary LLM (83.2--93.7\%). Second, candidate selection outperforms free generation for every LLM: accuracy rises from 8.1--50.7\% to 58.5--69.8\% for the open LLMs, and from 33.9--54.3\% to 83.2--93.7\% for the proprietary ones. Third, automatic scoring underrates free generation. It can count valid variants of an expansion, such as \heb{שקלים חדשים} for \heb{שקל חדש}, as wrong, and judging DictaLM's answers by hand raises its generation score from 32.8\% to 50.7\%.
